\documentclass[11pt,a4paper]{article}

\usepackage{../shared/cathedral-preamble}

\title{%
  From Zeta Zeros to Engineering:\\
  What the Cathedral's Mathematical Physics\\
  Means for Practicing Engineers%
}
\author{
  Jason Robert Gochanour
}
\date{April 20, 2026}

\begin{document}

\maketitle

\begin{abstract}
The Cathedral project establishes a machine-verified equivalence
between the Riemann Hypothesis and the decay of the Nyman--Beurling
distance, using techniques from spectral theory, harmonic analysis,
and variational principles. While the mathematical content is pure
number theory, the \emph{techniques} are drawn from the same toolkit
that practicing engineers use daily: Fourier analysis, eigenvalue
problems, variational methods, and structured matrix theory. This
paper maps the Cathedral's mathematical physics to five engineering
domains: signal processing, structural analysis, wireless
communications, control theory, and numerical methods. We do not
claim that the Cathedral's results have direct engineering
applications---they do not. Instead, we argue that the
\emph{mathematical structures} that appear in the Cathedral
illuminate engineering techniques from an unexpected angle, and
that engineers may find pedagogical and conceptual value in seeing
their daily tools applied to a 167-year-old unsolved problem.
\end{abstract}

\tableofcontents

% ================================================================
\section{Introduction: The Engineer's Toolkit in Disguise}
% ================================================================

Engineers solve problems with a canonical set of mathematical tools:

\begin{itemize}
  \item Fourier transforms and spectral analysis.
  \item Eigenvalue problems and matrix decompositions.
  \item Variational principles and optimization.
  \item Structured matrices (Toeplitz, Hankel, Gram).
  \item Digital filtering and windowing functions.
\end{itemize}

The Cathedral uses every one of these tools, applied to a problem
that has nothing to do with engineering---the distribution of prime
numbers. This paper makes the connection explicit, showing engineers
exactly where their tools appear in the proof architecture and what
the number-theoretic results look like through an engineering lens.

\subsection{What This Paper Is (and Isn't)}

\textbf{This paper is:} A translation guide that maps Cathedral
mathematics to engineering concepts, for engineers who are curious
about what a ``proof about primes'' looks like when you strip away
the number theory.

\textbf{This paper is not:} A claim that the Cathedral helps
engineers build better bridges, design better antennas, or write
better DSP code. The connection is structural (same tools) not
functional (same applications).

% ================================================================
\section{The Parseval Bridge as Signal Processing}
\label{sec:signal}
% ================================================================

\subsection{The Engineering Parseval Theorem}

Every electrical engineer learns Parseval's theorem in their first
signals course:
\[
  \int_{-\infty}^{\infty} |x(t)|^2\,dt =
  \int_{-\infty}^{\infty} |X(f)|^2\,df.
\]
The energy of a signal in the time domain equals the energy in the
frequency domain. This is the foundation of spectral analysis,
power spectral density, and filter design.

\subsection{The Cathedral's Version}

The Cathedral proves a specific instance:
\[
  \int_0^1 |1 - f_N(x)|^2\,dx =
  \frac{1}{2\pi} \int_{-\infty}^{\infty}
  |\hat{M}_N(1/2 + it)|^2\,dt,
\]
where $f_N(x)$ is a linear combination of fractional-part functions
$\{1/(kx)\}$, and $\hat{M}_N(s)$ is its Mellin transform on the
critical line.

\textbf{Engineering translation:}
\begin{center}
\begin{tabular}{@{}ll@{}}
\toprule
\textbf{Cathedral} & \textbf{Engineering} \\
\midrule
$\int_0^1 |1-f_N|^2\,dx$ & Time-domain error energy \\
$\hat{M}_N(1/2+it)$ & Frequency-domain spectrum \\
Mellin transform & Laplace-like transform ($x = e^{-u}$) \\
Critical line $\mathrm{Re}(s) = 1/2$ & Imaginary axis ($j\omega$) \\
$d_N^2 \to 0$ & Error energy $\to$ 0 \\
RH & The signal can be perfectly reconstructed \\
\bottomrule
\end{tabular}
\end{center}

The Mellin transform is related to the bilateral Laplace transform
by the substitution $x = e^{-u}$:
\[
  \mathcal{M}[f](s) = \int_0^\infty f(x)\,x^{s-1}\,dx
  = \int_{-\infty}^{\infty} f(e^{-u})\,e^{-su}\,du
  = \mathcal{L}[f(e^{-(\cdot)})](s).
\]
So the ``critical line'' $\mathrm{Re}(s) = 1/2$ in number theory
corresponds to the imaginary axis $s = 1/2 + j\omega$ in the
$s$-plane---exactly where control engineers evaluate transfer
functions for stability (the Nyquist criterion).

\subsection{The Window Function Connection}

The Cathedral's optimal test vectors $v_k$ are the Selberg sieve
weights:
\[
  v_k = \mu(k)\left(1 - \frac{\ln k}{\ln N}\right)^+,
\]
where $\mu(k)$ is the M\"obius function. The factor
$(1 - \ln k / \ln N)^+$ is a triangular window in log-space.

\textbf{Engineering translation:} This is a \emph{Bartlett window}
applied to the M\"obius spectrum. Every DSP engineer knows
Bartlett---it's the simplest window that tapers to zero at the
edges, reducing spectral leakage at the cost of widening the main
lobe.

The Cathedral's numerical experiments show this Bartlett window
achieves $Q_N \sim 12.45 \ln N$, while the optimal window achieves
$Q_N \sim 21.649 \ln N$. The ratio $12.45 / 21.649 \approx 0.575$
is analogous to the efficiency loss from using a Bartlett window
instead of a Dolph--Chebyshev or Kaiser window in DSP.

\textbf{Open question for engineers:} What is the optimal window
function $w(k)$ that minimizes $d_N^2$ for a given $N$? This is
equivalent to the optimal sidelobe suppression problem in antenna
array design.

% ================================================================
\section{The Gram Matrix as a Structural Problem}
\label{sec:structural}
% ================================================================

\subsection{Positive Definiteness and Stability}

The Cathedral proves that the augmented Gram matrix $H_N$ is
positive definite for all $N \geq 1$. In engineering:

\begin{itemize}
  \item \textbf{Structural engineering}: A positive definite
    stiffness matrix $K$ means the structure has no zero-energy
    deformation modes---it is stable.
  \item \textbf{Control theory}: A positive definite Lyapunov
    matrix $P$ certifies stability of a linear system $\dot{x} = Ax$.
  \item \textbf{Machine learning}: A positive definite kernel matrix
    $K$ means the kernel is valid (Mercer's theorem).
\end{itemize}

The Cathedral's proof that $H_N$ is PD uses the ``Factorial Nuke'':
the augmented entries involve $1/k!$, which decay so rapidly that
they dominate the spectral structure regardless of the off-diagonal
entries. This is standard in numerical analysis: diagonal dominance
implies positive definiteness.

\subsection{The Rayleigh Quotient}

The Rayleigh quotient $R(v) = v^T G v / v^T v$ gives the
minimum eigenvalue:
\[
  \lambda_{\min}(G) = \min_{v \neq 0} R(v).
\]
Engineers use this in:

\begin{itemize}
  \item \textbf{Vibration analysis}: The smallest eigenvalue of the
    generalized eigenvalue problem $Kv = \lambda Mv$ gives the
    fundamental natural frequency.
  \item \textbf{Buckling analysis}: The smallest eigenvalue of
    $K_g v = \lambda K v$ gives the critical buckling load.
  \item \textbf{Principal Component Analysis}: The eigenvalues of
    the covariance matrix determine variance explained.
\end{itemize}

The Cathedral's $\lambda_{\min}(G_N) \sim C / \ln N$ (where $C$ is
the quantum stiffness constant) means the Gram matrix becomes
increasingly ill-conditioned as $N$ grows---but never singular.
Engineers encounter exactly this behavior in mesh refinement for
finite element analysis: finer meshes produce larger, more
ill-conditioned stiffness matrices, but the physics ensures they
remain positive definite.

% ================================================================
\section{Spectral Gaps and Network Engineering}
\label{sec:networks}
% ================================================================

\subsection{Algebraic Connectivity}

In graph theory, the second-smallest eigenvalue of the Laplacian
matrix $L$ is the \emph{algebraic connectivity} (Fiedler value).
It measures how well-connected a network is:

\begin{itemize}
  \item $\lambda_2 = 0$: disconnected graph.
  \item $\lambda_2 > 0$: connected, with larger values meaning
    better connectivity.
  \item $\lambda_2$ controls mixing time of random walks on the graph.
\end{itemize}

The Cathedral's ``spectral gap'' $\lambda_{\min}(G_N) > 0$ for all
$N$ plays an analogous role: it ensures that the Gram matrix is
always invertible, meaning the best $L^2$ approximation always
exists and is unique.

\subsection{MIMO Channel Matrices}

In wireless communications, the channel capacity of a MIMO system
depends on the eigenvalues of $H H^*$, where $H$ is the channel
matrix. The eigenvalue distribution (Marchenko--Pastur for random
channels) determines how many independent data streams can be
multiplexed.

The Gram matrix $G_N$ is \emph{not} random---its entries are
deterministic functions of $\gcd(j,k)$. But its eigenvalue
distribution shares qualitative features with random matrix theory:
the bulk eigenvalues grow linearly while the minimum eigenvalue
grows logarithmically. This ``soft edge'' behavior is reminiscent
of the Tracy--Widom distribution at the edge of random matrix spectra.

% ================================================================
\section{Variational Methods: From RH to Finite Elements}
\label{sec:fem}
% ================================================================

\subsection{The Ritz Method}

The Cathedral finds the optimal test vector by minimizing:
\[
  d_N^2 = \min_{v \in \mathbb{R}^{N-1}}
  \int_0^1 \left(1 - \sum_{k=2}^N v_k \{1/(kx)\}\right)^2 dx
  = 1 - b^T G^{-1} b.
\]
This is a \emph{Ritz approximation}: minimize the energy functional
over a finite-dimensional subspace.

Engineers use exactly this structure in:

\begin{itemize}
  \item \textbf{Finite Element Method}: Minimize
    $\Pi(u) = \frac{1}{2} u^T K u - u^T f$ over the span of shape
    functions.
  \item \textbf{Method of Moments}: Project Maxwell's equations
    onto a finite basis to compute antenna patterns.
  \item \textbf{Galerkin methods}: Approximate PDE solutions by
    minimizing residuals in a Hilbert space.
\end{itemize}

The Cathedral's ``basis functions'' are $\{1/(kx)\}$ for
$k = 2, \ldots, N$. The ``target function'' is $\mathbf{1}_{(0,1)}$
(the constant function 1). The ``stiffness matrix'' is $G_N$. The
``load vector'' is $b_k = \int_0^1 \{1/(kx)\}\,dx$.

The entire Nyman--Beurling criterion reduces to: \emph{as the mesh
refines ($N \to \infty$), does the FEM solution converge to the
exact solution?} \RH{} says yes.

\subsection{Convergence Rate and Mesh Quality}

In FEM, the convergence rate depends on:
\begin{enumerate}
  \item The regularity of the target function (smoothness).
  \item The quality of the mesh (element aspect ratios).
  \item The condition number of the stiffness matrix.
\end{enumerate}

For the Cathedral:
\begin{enumerate}
  \item The target is $\mathbf{1}$---infinitely smooth.
  \item The ``mesh'' is the set of integers $\{2, \ldots, N\}$---with
    arithmetic structure instead of geometric structure.
  \item The condition number $\kappa(G_N) \sim N \ln N$---growing,
    but controlled.
\end{enumerate}

The convergence rate $d_N^2 \sim C / \ln N$ (under \RH) is
\emph{logarithmic}---much slower than polynomial FEM convergence
rates. This reflects the difficulty of the approximation problem:
the basis functions $\{1/(kx)\}$ are arithmetically structured, not
designed for optimal approximation.

% ================================================================
\section{Control Theory: Stability on the Critical Line}
\label{sec:control}
% ================================================================

\subsection{The Nyquist Connection}

In classical control theory, a system with transfer function $G(s)$
is stable if no poles lie in the right half-plane $\mathrm{Re}(s) > 0$.
The Nyquist stability criterion evaluates $G(j\omega)$ on the
imaginary axis to determine stability.

The Riemann Hypothesis asserts that the nontrivial zeros of
$\zeta(s)$ lie on the line $\mathrm{Re}(s) = 1/2$---the ``critical
line.'' Under the substitution $s \mapsto s + 1/2$, this becomes:
all zeros lie on the imaginary axis.

\textbf{Translation:} \RH{} says that a certain ``transfer function''
(related to $1/\zeta(s)$, the M\"obius generating function) has all
its zeros on the imaginary axis---the boundary of stability. The
primes are a marginally stable system.

\subsection{The Spectral Zeta Function as a Transfer Function}

Define $F(s) = \sum_{n=1}^{\infty} \mu(n) n^{-s}$. Then
$F(s) = 1/\zeta(s)$, and:
\begin{itemize}
  \item \RH{} $\iff$ all zeros of $F(s)$ have $\mathrm{Re}(s) = 1/2$.
  \item $F(s)$ converges absolutely for $\mathrm{Re}(s) > 1$ (the
    ``high-frequency'' region).
  \item The behavior on $1/2 < \mathrm{Re}(s) < 1$ is the ``critical
    strip''---the frequency band where the system's behavior is
    uncertain.
\end{itemize}

This is reminiscent of robust control: we know the system is stable
at high frequencies, but the critical behavior is in a finite
bandwidth around the stability boundary.

% ================================================================
\section{Numerical Methods: Condition Numbers and the Quantum Stiffness}
\label{sec:numerical}
% ================================================================

\subsection{The Condition Number Problem}

Numerical engineers care about the condition number
$\kappa(A) = \|A\| \cdot \|A^{-1}\|$, which controls the
amplification of roundup errors in solving $Ax = b$.

The Cathedral's Gram matrix has condition number
$\kappa(G_N) \sim N \ln N / C$, where $C \approx 0.046$ is the
quantum stiffness constant. This is a moderately ill-conditioned
system: worse than a Hilbert matrix ($\kappa \sim e^{3.5N}$) but
much better than the Vandermonde matrix ($\kappa \sim N^N$).

\textbf{Engineering implication:} The Gram matrix can be solved
numerically in double precision for $N \lesssim 10^{12}$ before
roundoff destroys the solution. This is why the Cathedral's Rust
experiments use exact rational arithmetic---to avoid conditioning
issues entirely.

\subsection{Iterative Methods}

The Gram matrix is symmetric positive definite, which makes it
ideal for the Conjugate Gradient method. The convergence rate
depends on $\sqrt{\kappa}$:
\[
  \|x_k - x_*\| \leq 2 \left(\frac{\sqrt{\kappa} - 1}
  {\sqrt{\kappa} + 1}\right)^k \|x_0 - x_*\|.
\]
For $\kappa \sim N \ln N$, this gives convergence in
$O(\sqrt{N \ln N})$ iterations---manageable for moderate $N$.

% ================================================================
\section{The Dictionary}
\label{sec:dictionary}
% ================================================================

\begin{center}
\small
\begin{tabular}{@{}lll@{}}
\toprule
\textbf{Cathedral} & \textbf{Engineering Domain} & \textbf{Engineering Concept} \\
\midrule
Parseval Bridge & Signal processing & Time-frequency duality \\
$d_N^2 \to 0$ & DSP & Error energy convergence \\
Selberg weights & DSP & Bartlett window function \\
Gram matrix $G_N$ & Structural / FEM & Stiffness matrix \\
$\lambda_{\min}(G_N)$ & Structural / FEM & Fundamental frequency \\
$H_N$ positive definite & Control & Lyapunov stability \\
Rayleigh quotient & Vibrations & Natural frequency bound \\
Critical line $\mathrm{Re}(s)=1/2$ & Control & Nyquist imaginary axis \\
$1/\zeta(s)$ & Control & Transfer function \\
RH & Control & Marginal stability \\
$\kappa(G_N)$ & Numerical methods & Condition number \\
$Q_N / \ln N \to C$ & Materials & Spectral stiffness constant \\
Ritz approximation & FEM & Galerkin projection \\
Axiom elimination & V\&V & Assumption retirement \\
\bottomrule
\end{tabular}
\end{center}

% ================================================================
\section{What Engineers Cannot Learn from the Cathedral}
\label{sec:cannot}
% ================================================================

For intellectual honesty:

\begin{enumerate}
  \item The Cathedral does not provide new signal processing
    algorithms. The Parseval Bridge is a specific application of a
    known theorem.
  \item The Gram matrix structure (entries involving $\gcd$ and
    $\operatorname{lcm}$) does not appear in typical engineering
    problems. No bridge or antenna produces a Gram matrix.
  \item The convergence rate $d_N^2 \sim C / \ln N$ is far too slow
    for any engineering approximation scheme. Engineers would never
    accept logarithmic convergence.
  \item The ``quantum stiffness constant'' $C \approx 0.046$ is a
    number-theoretic constant with no known physical realization.
\end{enumerate}

The value for engineers is \emph{conceptual}, not \emph{practical}:
seeing familiar tools applied to an unfamiliar problem can deepen
understanding of why those tools work and what their fundamental
limitations are.

% ================================================================
\section{Conclusion}
% ================================================================

The Cathedral is pure mathematics, verified by machine. But its
tools are the engineer's tools: Fourier analysis, eigenvalue
problems, variational methods, condition numbers, and window
functions. The primes are not a bridge or a circuit or an antenna,
but they respond to the same mathematical forces.

For the practicing engineer, the Cathedral offers a mirror:
\emph{the mathematics you use to design filters and analyze
structures is the same mathematics that governs the distribution of
prime numbers.} This is not a coincidence. It is a reflection of the
universality of harmonic analysis---the deepest tools work
everywhere, from the critical line of the zeta function to the
critical frequency of a low-pass filter.

\begin{quote}
\emph{The engineer's toolkit was built for bridges.\\
It turns out the primes are also a bridge---\\
between order and chaos, between structure and randomness.\\
The tools work because the bridge is real.}
\end{quote}

\vspace{0.5cm}

\begin{thebibliography}{99}

\bibitem{oppenheim}
A.V.~Oppenheim and R.W.~Schafer, \emph{Discrete-Time Signal
Processing}, 3rd ed., Pearson, 2009.

\bibitem{strang}
G.~Strang, \emph{Linear Algebra and Its Applications}, 4th ed.,
Cengage, 2006.

\bibitem{hughes}
T.J.R.~Hughes, \emph{The Finite Element Method: Linear Static and
Dynamic Finite Element Analysis}, Dover, 2000.

\bibitem{trefethen}
L.N.~Trefethen and D.~Bau, \emph{Numerical Linear Algebra}, SIAM,
1997.

\bibitem{ogata}
K.~Ogata, \emph{Modern Control Engineering}, 5th ed., Pearson, 2010.

\bibitem{cathedral-physics}
J.R.~Gochanour, \emph{The Cathedral: A Dictionary Between Pure
Mathematics and Quantum Field Theory}, 2026.

\end{thebibliography}

\end{document}
